% 本文件由 LaTeX 排版 Agent 根据有效 translation.json 自动生成。
% author=Moses of Chorene (c. 400-c. 490); work_id=0859; title=亚美尼亚史

\chapter{\WorkTitle{亚美尼亚史}}

% structure: p0001 source=h1 target=omitted reason=page-title-already-rendered-by-parent

\Emph{\Person{科伦的摩西}}。\par

% structure: p0003 source=h2 target=section reason=relative-heading-rank; base=h2

\section{一、\Person{阿布加尔}在位时期；亚美尼亚完全沦为罗马属国；与\Person{黑落德}军队的战争；其兄之子\Person{若瑟}被杀。}

\Person{阿布加尔}，\Person{阿尔卡姆}之子，在波斯王\Person{阿尔查维尔}在位第二十年登基。这位\Person{阿布加尔}因其极大的温和与智慧，并因其身材高大，被称为\OriginalTerm{阿瓦克-艾尔}{Avak-air}（伟人）。希腊人和叙利亚人发音不准，称他为\Emph{\Person{阿布加尔}}。他在位第二年，亚美尼亚各地均成为罗马属国。奥古斯都皇帝下令，正如我们在\BibleBook{路加}{福音}中所读到的，要登记各处所有百姓。为此派往亚美尼亚的罗马专员携带奥古斯都皇帝的雕像，并将其立于所有庙宇中。就在此时，我们的救主耶稣基督、天主之子降生成人。\par

同一时期，\Person{阿布加尔}与\Person{黑落德}之间发生冲突：因为\Person{黑落德}希望他的雕像能立在亚美尼亚庙宇中凯撒雕像旁边。\Person{阿布加尔}反对此要求。而且，\Person{黑落德}只是在寻找攻击\Person{阿布加尔}的借口：他派遣了一支由色雷斯人和日耳曼人组成的军队入侵波斯领土，并命令他们穿过\Person{阿布加尔}的领地。但\Person{阿布加尔}并未屈服，而是抵抗，声称皇帝的命令是要军队穿越沙漠进入波斯。\Person{黑落德}愤慨不已，又因自己对基督的恶行而陷入困境，正如\Person{约瑟夫}所述，他无法亲自行动，便派他的侄子（他曾将女儿许配给他，此女最初嫁给了他的兄弟\Person{费若尔}）前往。\Person{黑落德}的副将率领大军急速抵达美索不达米亚，在\Place{普克南}省的营地与\Person{阿布加尔}相遇，战死沙场，其军队溃逃。不久，\Person{黑落德}死去；其子\Person{阿尔刻劳}被\Person{奥古斯都}任命为犹太地总督。\par

% structure: p0006 source=h2 target=section reason=relative-heading-rank; base=h2

\section{二、\Place{厄德萨}城的建立；光照者家族简述。}

不久之后，\Person{奥古斯都}去世，\Person{提比略}继任罗马皇帝。\Person{日耳曼尼库斯}成为凯撒，拖着\Person{阿尔查维尔}和\Person{阿布加尔}王国的王子们，为他们之间的战争举行凯旋式，这些王子曾杀死\Person{黑落德}的侄子。\Person{阿布加尔}愤慨，计划叛乱并准备战斗。他在亚美尼亚警戒军队驻扎的土地上建造了一座城，此前\Person{卡西乌斯}曾在此处防御幼发拉底河：这座新城被称为\Place{厄德萨}。\Person{阿布加尔}将他的宫廷从\Place{梅兹平}迁至此地，并带去了所有神祇：\Person{纳博克}、\Person{贝耳}、\Person{帕特尼卡赫}和\Person{塔拉塔}，以及附属于庙宇的学校书籍，甚至王室档案。\par

此后，\Person{阿尔查维尔}去世，其子\Person{阿尔达凯斯}统治波斯人。虽然从时间顺序来看，这并不符合历史的顺序，甚至也不符合我们开始编年史的顺序，但既然我们正在论述\Person{阿尔查维尔}王的子孙，即他儿子\Person{阿尔达凯斯}的血脉，为了向这些王子表示敬意，我们将提前时间，将他们置于\Person{阿尔达凯斯}附近，以便读者知晓他们属于同一血脉，即勇敢的\Person{阿尔查格}的血脉；然后，我们将指出他们的祖先来到亚美尼亚的时间，即\Person{加伦尼}家族和\Person{苏伦尼}家族，\Person{圣额我略}和\Person{加姆萨里}家族由此而出，当我们按照事件顺序来到他们出现的国王统治时期时，再行叙述。\par

\Person{阿布加尔}的叛乱计划并未成功；因为，波斯王国中他的亲属之间发生了骚乱，他便率军前往平息并结束纷争。\par

% structure: p0010 source=h2 target=section reason=relative-heading-rank; base=h2

\section{三、\Person{阿布加尔}来到东方，维护\Person{阿尔达凯斯}在波斯王位；使他的兄弟们和好，我们的光照者及其亲属由此而出。}

\Person{阿布加尔}到了东方，发现波斯王位上坐着\Person{阿尔查维尔}之子\Person{阿尔达凯斯}，而\Person{阿尔达凯斯}的兄弟们正在与他争战：因为这位王子想在他的后代中统治他们，而他们不同意。因此，\Person{阿尔达凯斯}从四面八方包围他们，将死亡的利剑悬在他们头上；他们的军队与他们的其他亲属和盟友之间发生了分裂和冲突：因为\Person{阿尔查维尔}王有三个儿子和一个女儿；其中长子就是\Person{阿尔达凯斯}王本人，次子\Person{加伦}，三子\Person{苏伦}；他们的妹妹名叫\Person{戈赫姆}，是所有阿里克人的将军之妻，这位将军是由他们的父亲\Person{阿尔查维尔}选定的。\par

\Person{阿布加尔}说服\Person{阿尔查维尔}的儿子们讲和；他在他们之间安排了条件和约定：\Person{阿尔达凯斯}将按他所提议的，与他的后代一同统治，而他的兄弟们将被称为\OriginalTerm{巴赫拉夫}{Bahlav}，取自他们的城镇和广阔肥沃的土地之名，以使他们的总督区成为波斯所有总督区中第一等、地位最高者，因为是真正的王室血脉。条约和誓言规定：如果\Person{阿尔达凯斯}没有男嗣，他的兄弟们应继承王位；在\Person{阿尔达凯斯}的统治家族之后，他的兄弟们分为三支，命名如下：\Person{加伦·巴赫拉夫}家族、\Person{苏伦·巴赫拉夫}家族，以及他们妹妹的家族——\Person{阿斯巴哈比德·巴赫拉夫}家族，这一家族以其丈夫的领地之名命名。\par

\Person{圣额我略}据说出自\Person{苏伦·巴赫拉夫}家族，而\Person{加姆萨里}家族出自\Person{加伦·巴赫拉夫}家族。我们将在后文叙述这些人物到来时的情形，此处仅提及他们的名字与\Person{阿尔达凯斯}相关联，以便你知道这些伟大的家族确实是\Person{瓦加尔沙格}的血脉，即伟大的\Person{阿尔查格}的后裔，\Person{瓦加尔沙格}的兄弟。\par

一切安排妥当后，\Person{阿布加尔}带着条约文书返回自己的领地；身体并非完全健康，而是遭受着严重的痛苦。\par

% structure: p0015 source=h2 target=section reason=relative-heading-rank; base=h2

\section{四、\Person{阿布加尔}从东方返回；他在一场战争中援助\Person{阿雷塔斯}对抗分封侯\Person{黑落德}。}

\Person{阿布加尔}从东方回来时，得知罗马人怀疑他去那里是为招募军队。于是，他向罗马特使说明了他前往波斯的缘由，以及\Person{阿尔达赫斯}与其兄弟之间订立的条约；但他的话未被采信，因为他的敌人\Person{比拉多}、分封侯\Person{黑落德}、\Person{吕撒尼雅}和\Person{斐理伯}控告他。\Person{阿布加尔}回到自己的城市\Place{厄德萨}后，与\Place{培特拉}王\Person{阿勒塔斯}结盟，并派出一支由\Person{科斯兰·阿尔兹鲁尼}指挥的辅助部队攻打\Person{黑落德}。\Person{黑落德}最初娶了\Person{阿勒塔斯}的女儿，后来休了她，又娶了\Person{黑落狄亚}，当时她的丈夫还在世，此事导致他处死了\Person{洗者若翰}。因此，\Person{黑落德}与\Person{阿勒塔斯}之间因\Person{阿勒塔斯}女儿所受的冤屈而爆发战争。\Person{黑落德}的军队遭到猛攻而溃败，得益于勇敢的亚美尼亚人的帮助；仿佛出于天主眷顾，为\Person{洗者若翰}的死报了仇。\par

% structure: p0017 source=h2 target=section reason=relative-heading-rank; base=h2

\section{五、\Person{阿布加尔}派王子往\Person{马里努斯}处；这些使节见到了我们的救主基督；\Person{阿布加尔}归化之始。}

此时，\Person{斯托罗格}之子\Person{马里努斯}被皇帝任命为腓尼基、巴勒斯坦、叙利亚和美索不达米亚的总督。\Person{阿布加尔}派了他的两位主要官员——\Person{阿赫特兹尼克}亲王\Person{马里哈普}、\Person{阿巴胡尼}家族首领\Person{查姆查克拉姆}，以及他的亲信\Person{阿南}——前往\Place{培特库皮内}城，向\Person{马里努斯}说明\Person{阿布加尔}前往东方的缘由，向他出示\Person{阿尔达赫斯}与其兄弟订立的条约，并同时请求\Person{马里努斯}的支持。使节们在\Place{厄娄特罗}城找到了这位罗马总督；他友好而隆重地接待了他们，并给\Person{阿布加尔}如下答复：\Quote{只要你们注意按时缴纳贡赋，在这件事上就不必惧怕皇帝。}\par

回程途中，亚美尼亚使节们前往\Place{耶路撒冷}，想亲眼看看我们的救主基督，因他们被祂的奇迹传闻所吸引。他们亲眼目睹了这些奇事，便将其告知了\Person{阿布加尔}。\Person{阿布加尔}满怀崇敬，真诚地相信\Person{耶稣}确实是天主子，说道：\Quote{这些奇事并非人力所为，而是天主之能。不，世人中无人能复活死者：唯天主有此大能。}\Person{阿布加尔}全身感到剧烈的疼痛，这病是在\Place{波斯}患上的，已有七年多；从人那里他未得到任何医治。于是\Person{阿布加尔}写了一封恳求的信给\Person{耶稣}，祈求祂来治愈他的病痛。以下是这封信——\par

% structure: p0020 source=h2 target=section reason=relative-heading-rank; base=h2

\section{六、\Person{阿布加尔}致救主\Person{耶稣}基督的书信}

\Person{阿尔卡姆}之子、此地王公\Person{阿布加尔}，致在\Place{耶路撒冷}地区显现的人类的救主和恩主\Person{耶稣}，致敬——\par

\Quote{我听到关于祢的事，以及祢手所行的医治，不用药物、不用草药：因为据说，祢使盲者复明，跛者行走，癞者洁净；祢驱逐不洁之神，治愈长期受顽疾折磨的可怜人；祢甚至复活死人。我既听到祢行的这一切奇事，便断定：祢或是自天降下的天主，行此大事；或是天主子，行此奇迹。因此我写信给祢，求祢屈尊到我这里来，治愈我的病痛。我也听说犹太人对祢不满，想要把祢交出受折磨。我有一座城，虽小却宜人，足够我们二人居住。}\par

送信的使者在\Place{耶路撒冷}遇见了\Person{耶稣}，福音中的话证实了此事：\Quote{有些外邦人来见\Person{耶稣}，但听见的人不敢将所听见的告诉\Person{耶稣}，就告诉了\Person{斐理伯}和\Person{安德肋}，他们两人便将一切禀报了他们的师傅。}\par

救主当时并未接受邀请，但祂认为应赐予\Person{阿布加尔}一封回信，内容如下：——\par

% structure: p0025 source=h2 target=section reason=relative-heading-rank; base=h2

\section{七、对\Person{阿布加尔}书信的答复，此信由宗徒\Person{多默}奉救主之命写给这位王公。}

\Quote{那没有看见我而相信的，是有福的！因为经上论到我记载说：\InnerQuote{看见我的人不会信我，而没有看见我的人却要信而生活。}至于你所写请我到你那里去的事，我必须在此完成我被派遣所做的一切；当我完成这一切后，我将升到那派遣我者那里去；当我离去后，我会派我的一位门徒来，他将治愈你的疾病，并使你和与你同在的人都得生命。}\Person{阿南}，\Person{阿布加尔}的信使，将这封信以及救主的肖像带给\Person{阿布加尔}，这幅画像至今仍在\Place{厄德萨}城中。\par

% structure: p0027 source=h2 target=section reason=relative-heading-rank; base=h2

\section{八、宗徒\Person{达陡}在\Place{厄德萨}的宣讲；五封书信的副本。}

我们的救主升天以后，十二位宗徒之一的\Person{多默}宗徒，按照主的圣言，从七十六位门徒中派遣了\Person{达陡}前往\Place{厄德撒}城，为治愈\Person{阿布加尔}并宣讲福音。\Person{达陡}来到一位名叫\Person{托彼特}的犹太亲王的家中，据说此人属于\OriginalTerm{帕克拉杜尼}{Pacradouni}家族。\Person{托彼特}离开\Place{阿尔卡姆}后，并未像其余亲属那样背弃犹太教，而是一直遵行其律法，直到他相信基督的时候。不久，\Person{达陡}的名声传遍全城。\Person{阿布加尔}得知他的到来，便说：\Quote{这的确就是耶稣写信给我提到的那一位；}于是\Person{阿布加尔}立即派人去请宗徒。当\Person{达陡}进入时，\Person{阿布加尔}眼前呈现出宗徒面容上的奇妙景象；国王从宝座上站起，面伏于地，在\Person{达陡}面前俯伏叩拜。这一景象令在场所有亲王大为惊奇，因为他们并不知道这异象的事。\Person{阿布加尔}对\Person{达陡}说：\Quote{你真的是那永远可称颂的耶稣的门徒吗？你就是祂应许要派到我这里来的那位，并能够医治我的疾病吗？} \Quote{是的，}\Person{达陡}回答说，\Quote{你若相信\Person{耶稣基督}、天主子，你心中的愿望就必得成就。} \Quote{我已经相信了耶稣，}\Person{阿布加尔}说，\Quote{我已经相信了圣父；因此我本想率领我的军队去消灭那些钉死耶稣的犹太人，若不是因为罗马人的势力拦阻了我。}\par

此后，\Person{达陡}开始向国王及其全城宣讲福音；他按手在\Person{阿布加尔}身上，治好了他；也治好了一位患痛风的人，即\Person{阿布杜}亲王，此人在全朝中备受尊敬。他还治好了城中所有患病和软弱的人，所有人都相信了\Person{耶稣基督}。\Person{阿布加尔}受了洗，全城的人也与他一同受了洗，假神的庙宇被关闭，所有放置在祭台和柱子上的偶像雕像都被芦苇遮盖而隐藏起来。\Person{阿布加尔}并未强迫任何人接受信仰，然而信徒的人数日日增多。\par

宗徒\Person{达陡}为一位名叫\Person{阿塔乌斯}的丝制头巾制造商施洗，祝圣了他，任命他在\Place{厄德撒}\Emph{服务}，并把他留在国王那里代替自己。\Person{达陡}从\Person{阿布加尔}那里获得了特许状——国王希望所有人都听从基督的福音——之后，便前去找\Person{阿布加尔}的姐妹之子\Person{萨纳特鲁格}，这位亲王已被任命管辖该地区和军队。\Person{阿布加尔}还高兴地给罗马皇帝\Person{提庇留}写了一封如下内容的信：——\par

% structure: p0031 source=h3 target=subsection reason=relative-heading-rank; base=h2

\subsection{阿布加尔致提庇留的信}

\Person{阿布加尔}，\Place{亚美尼亚}王，致我主\Person{提庇留}，\Place{罗马}皇帝，安好：——\par

\Quote{我知道陛下无所不知，但我作为您的朋友，愿以书信使您更清楚地了解事实。居住在\Place{巴勒斯坦}各地的犹太人已经把\Person{耶稣}钉在十字架上：耶稣是无罪的，耶稣在行了如此多的善事、如此多的奇事和奇迹造福他们，甚至使死人复活之后。请确信，这些并非一个凡人的能力所及，而是出于天主。在他们钉死祂的时候，太阳变暗了，大地震动、摇撼；耶稣自己，三天之后，从死者中复活，并显现给许多人。如今，到处都有祂的名号——由祂的门徒呼求——产生最大的奇迹：我个人所经历的事就是最明显的证明。陛下从此知道对于犯下这罪的犹太民族将来应当如何处理；陛下也知道是否应当向普天下颁布命令，要敬拜基督为真天主。祝平安康泰。}\par

% structure: p0034 source=h3 target=subsection reason=relative-heading-rank; base=h2

\subsection{提庇留对阿布加尔来信的答复}

\Person{提庇留}，\Place{罗马}皇帝，致\Person{阿布加尔}，\Place{亚美尼亚}王，安好：——\par

\Quote{你的亲切来信我已阅毕，我愿向你表示谢意。虽然我们已从数人那里听到这些事实，但\Person{比拉多}已将\Person{耶稣}的奇迹正式告知我们。他向我们证实，耶稣从死者中复活后，被许多人承认为天主。因此，我自己也愿做你所建议的事；然而，由于\Place{罗马}人的惯例是不凭君主的命令接纳神祇，而只有在元老院全体讨论和审查之后才予接纳，我便将此事提交元老院，但他们轻蔑地否决了，无疑是因为此事未先经他们审议。但是，我们已命令所有认可\Person{耶稣}的人，把祂列入诸神之中。我们威胁任何胆敢毁谤基督徒的人，将要处死。至于那个胆敢把\Person{耶稣}钉死的犹太民族——据我所闻，祂非但不该受十字架和死亡，反而应受尊荣、应受人的朝拜——等我从与反叛的\Place{西班牙}的战争中脱身之后，我将查办此事，并按照犹太人应得的来处置他们。}\par

% structure: p0037 source=h3 target=subsection reason=relative-heading-rank; base=h2

\subsection{阿布加尔另致提庇留的信}

\Person{阿布加尔}，\Place{亚美尼亚}王，致我主\Person{提庇留}，\Place{罗马}皇帝，安好：——\par

\Quote{我已收到陛下写来的信，并赞赏您智慧所发出的命令。如果您不生气，我要说元老院的行为极其可笑荒谬：因为按照元老们的意见，神性是要经过审查和人的投票才能赋予的。那么，如果天主不合人的意，祂就不能是天主，因为天主是要由人来审判和证明的。我想，我的主上会认为应该另派一位总督到\Place{耶路撒冷}代替\Person{比拉多}，他应当被耻辱地赶出您所赐予他的权位；因为他顺从了犹太人的意愿，未经您的命令就不义地把基督钉死在十字架上。愿您健康，这是我的愿望。}\par

\Person{阿布加尔}写了这封信后，将一份信函连同其他信函的副本存入他的档案。他还写信给在\Place{巴比伦}的\Place{亚述}王年轻的\Person{尼尔塞}如下：——\par

% structure: p0041 source=h3 target=subsection reason=relative-heading-rank; base=h2

\subsection{阿布加尔致尼尔塞的信}

\Person{阿布加尔}，\Place{亚美尼亚}王，致我子\Person{尼尔塞}，安好：——\par

\Quote{我已收到你的信函和致谢。我释放了贝若泽，免除了他的罪过：若你满意，可将尼尼微的统治权交予他。至于你信中所提，请你将那能行奇迹、并宣讲另一位超越水火之神者派来，好使你亲眼见他、亲耳听他，我对你说：他并非按照人间医术的医生，他是创造水火的天主之子的一位门徒；他被派遣至亚美尼亚地区。但他的一位主要同伴，名叫西满，被派往波斯地区。寻找他吧，你将听到他的宣讲，你及你的父亲阿尔达赫斯皆然。他必治愈你所有疾病，并为你指示生命之道。}\par

% structure: p0044 source=h3 target=subsection reason=relative-heading-rank; base=h2

\subsection{阿布加尔致阿尔达赫斯的信}

阿布加尔又致书波斯王阿尔达赫斯说：——\par

亚美尼亚王阿布加尔，致我兄波斯王阿尔达赫斯，请安：——\par

\Quote{我知道你已听闻耶稣基督、天主之子，祂被犹太人钉在十字架上，却从死者中复活，并派遣祂的门徒遍及天下训导万民。祂的一位首席门徒，名叫西满，正在陛下的疆域内。寻找他吧，你必寻见，他必治愈你所有疾病，并为你指示生命之道，你、你的兄弟及所有乐意服从你的人，都将相信他的话。想到我的骨肉之亲也将成为我精神上的亲属和朋友，我深感欣慰。}\par

阿布加尔尚未收到这些信的回音便已驾崩，在位凡三十八年。\par

% structure: p0049 source=h2 target=section reason=relative-heading-rank; base=h2

\section{九、我等宗徒的殉道}

阿布加尔死后，亚美尼亚王国分为二：阿布加尔之子阿纳农统治埃德萨，其外甥萨纳特鲁格统治亚美尼亚。他们时代所发生的事，前人已有所记述：宗徒抵达亚美尼亚，萨纳特鲁格归信及因畏惧亚美尼亚总督而背道，宗徒及其同伴在查瓦尔昌地区（今称阿尔达兹）殉道，岩石裂开接纳宗徒遗体，门徒移走遗体，葬于平原，王之女桑图赫德在路旁殉道，两位圣人的遗骸显现并移至岩石——凡此种种，均如我们所说，早在我们之前已由他人记述；我们认为不必在此重述。同样，阿布加尔之子下令将宗徒的门徒阿泰奥斯在埃德萨处死的记述，前人亦已述及。\par

那位在其父死后继位的王子，并未继承其父的德行：他重开偶像庙宇，信奉异教。他派人告诉阿泰奥斯：\Quote{请用金线交织的布料为我制作一顶头饰，就如同你从前为我父亲制作的那样。}阿泰奥斯回答说：\Quote{我的手不会为一位不配的王子制作头饰，因他不敬拜永生的天主基督。}\par

国王随即命一武士斩断阿泰奥斯的双脚。武士前往，见那圣者坐在教师的座椅上，便用剑砍断其双腿，圣人当即气绝。我们在此略述，此事亦已久为他人所记。随后宗徒巴尔多禄茂来到亚美尼亚，在阿雷潘城殉道。至于被派往波斯的西满，我无法确述其所行及何处殉道。有说一位宗徒西满在韦里奥斯波尔殉道。此说是否确实？圣人为何来此？我不得而知；仅略及此，俾知我尽其所能告知一切必要之事。\par

% structure: p0053 source=h2 target=section reason=relative-heading-rank; base=h2

\section{十、萨纳特鲁格统治；阿布加尔子女被杀害；赫肋纳公主}

萨纳特鲁格即位后，借助曾推举他的勇武的帕克拉杜尼和阿尔兹鲁尼家族招募军队，前去征讨阿布加尔的子女，欲将整个王国据为己有。正当萨纳特鲁格忙于此事时，仿佛是上智安排的结果，阿泰奥斯之死得到了报复：阿布加尔之子正在埃德萨王宫顶部竖立一根大理石柱，他在柱下指挥时，柱子从工匠手中脱落，砸在他身上，碾碎了他的双脚。\par

城中居民随即遣使向萨纳特鲁格求和，请他承诺不干扰他们信奉基督宗教的自由，以此作为条件，他们将献出城池及王室的财宝。萨纳特鲁格应允，但最终背弃誓言。萨纳特鲁格将阿布加尔家族所有男丁悉数斩于剑下，只留下女儿们，他将她们接出城，安置在哈赫迪扬克地区。至于阿布加尔的原配，名赫肋纳，他将她送往哈兰城，并将整个美索不达米亚的统治权留给她，以纪念他通过赫肋纳从阿布加尔所受的恩惠。\par

赫肋纳虔敬如同其夫阿布加尔，不愿居于拜偶像者中间；在克劳狄在位时期，恰逢阿加波所预言的饥荒，她前往耶路撒冷，用其全部财宝在埃及购买大量谷物，分给穷人。此事有若瑟夫为证。赫肋纳的坟墓，确实引人注目，至今仍可见于耶路撒冷城门前。\par

% structure: p0057 source=h2 target=section reason=relative-heading-rank; base=h2

\section{十一、梅兹皮内城重建；萨纳特鲁格之名；其死亡}

在萨纳特鲁格的所有事迹中，我们认为唯有建筑梅兹皮内城值得记念；因该城遭地震摇撼，萨纳特鲁格将其拆毁，更加壮丽地重建，并围以双层城墙与壁垒。萨纳特鲁格命人在城中心竖立其雕像，手中握一枚钱币，以示：\Quote{我所有财宝已用于建城，仅剩此一枚钱币留于己。}\par

但是，为什么这位王子被称为\Person{萨纳特鲁格}？我们要告诉你：因为\Person{阿布加尔}的妹妹\Person{奥泰亚}，在冬季旅行于亚美尼亚时，在\Place{戈尔图克}山遭到一阵雪旋风；暴风雪使他们彼此失散，无人知道同伴被吹到了哪里。王子的乳母\Person{萨诺德}，\Person{皮乌拉德·帕克拉杜尼}之妹，\Person{霍斯兰·阿尔兹鲁尼}之妻，将王室婴儿（因为\Person{萨纳特鲁格}尚在摇篮中）抱在怀中，在雪下与他一起待了三天三夜。传说抓住了这一情况：它讲述了一种动物，一种新的物种，极其洁白，由众神派来，守护着孩子。但据我们所知，事实是这样的：一只在搜寻人员中的白狗，找到了孩子和他的乳母；因此王子被称为\Person{萨纳特鲁格}，这个名字取自他乳母的名字（以及亚美尼亚语名称\Emph{杜尔克}，意为礼物），仿佛意指萨诺德的礼物。\par

\Person{萨纳特鲁格}，在波斯王\Person{阿尔达切斯}十二年登基，在位三十年后，在狩猎时被一支箭射穿肠子而死，仿佛是为他对圣洁的女儿所施加的折磨而受罚。缮写员\Person{阿普沙塔尔}之子\Person{盖鲁普纳}，收集了所有这些发生在\Person{阿布加尔}和\Person{萨纳特鲁格}时期的事实，并将它们存入了\Place{厄德萨}的档案中。\par

\SourceRecord{Translated by B.P. Pratten.；Ante-Nicene Fathers；Vol. 8.；Edited by Alexander Roberts, James Donaldson, and A. Cleveland Coxe.；Buffalo, NY: Christian Literature Publishing Co.,；1886.；<http://www.newadvent.org/fathers/0859.htm>.}
